% from ega4-11-fr.tex
\oldpage[IV-3]{128}

soit surjectif~; cela signifie aussi que, pour tout entier $n$,
l'homomorphisme canonique de $k(\mathfrak q)$-modules
\[
\left(\bigoplus_{-N\leq i\leq N}(M_i\otimes_{B_0}B_{n-i})\right)
\otimes_Ak(\mathfrak q)\longrightarrow M_n\otimes_Ak(\mathfrak q)
\]
est surjectif. Or, $M_n$ est un $B_0$-module de type fini, $B_0$ un anneau
local et $\mathfrak qB_0\subset\mathfrak m$~; donc le lemme de Nakayama prouve
que chacun des homomorphismes canoniques
\[
\bigoplus_{-N\leq i\leq N}(M_i\otimes_{B_0}B_{n-i})\longrightarrow M_n
\]
est surjectif, d'où la conclusion.

\begin{corollary}[11.2.9.4]
\label{IV.11.2.9.4-fr}
Sous les hypothèses de (11.2.9.3), on suppose en outre que chacun des $B_i$ et
chacun des $M_i$ soit un $B_0$-module de présentation finie, et que $M$ soit
un $A$-module plat. Pour que $M$ soit un $B$-module de présentation finie, il
faut et il suffit que $M\otimes_Ak(\mathfrak q)$ soit un
$(B\otimes_Ak(\mathfrak q))$-module de présentation finie.
\end{corollary}

En vertu de (11.2.9.3), si la condition de l'énoncé est vérifiée, $M$ est un
$B$-module de type fini, donc il existe un $B$-module libre \emph{gradué} de
type fini $L$ (ayant donc une base finie formée d'éléments homogènes) et un
homomorphisme gradué surjectif de degré 0, $u:L\to M$. Soit $R=\ker(u)$, qui
est un $B$-module gradué~; il y a donc un nombre fini d'entiers $m_j$
($1\leq j\leq r$) tels que pour chaque entier $i$, $R_i$ soit le noyau d'un
homomorphisme surjectif
\[
\bigoplus_{1\leq j\leq r}B_{i+m_j}\longrightarrow M_i~;
\]
on conclut alors de l'hypothèse sur les $M_i$ et les $B_i$ que $R_i$ est un
$B_0$-module de type fini (Bourbaki, \emph{Alg. comm.}, chap.~I, \S~2,
n$^\circ$~8, lemme 9). Pour prouver que $R$ est un $B$-module de type fini,
notons qu'en vertu de l'hypothèse de platitude et de
($\mathbf 0_{\mathrm I}$, 6.1.2) la suite
\[
0\longrightarrow R\otimes_Ak(\mathfrak q)
\longrightarrow L\otimes_Ak(\mathfrak q)
\longrightarrow M\otimes_Ak(\mathfrak q)\longrightarrow0
\]
est exacte, et l'hypothèse entraîne donc (Bourbaki, \emph{loc. cit.}) que
$R\otimes_Ak(\mathfrak q)$ est un $(B\otimes_Ak(\mathfrak q))$-module de type
fini~; il suffit donc d'appliquer à $R$ le lemme (11.2.9.3).

\medskip
\noindent\textbf{III)} \emph{Réduction au cas où les homomorphismes de
transitions $\varphi_{\mu\lambda}:A_\lambda\to A_\mu$ (pour
$\lambda\leq\mu$) sont injectifs.} --- Soit $A'_\lambda$ l'image de
$A_\lambda$ par l'homomorphisme canonique $\varphi_\lambda:A_\lambda\to A$~;
il est clair que les $A'_\lambda$ forment un système inductif de sous-anneaux
noethériens de $A$, dont la limite inductive est $A$, et les homomorphismes de
transition $A'_\lambda\to A'_\mu$ (pour $\lambda\leq\mu$) sont injectifs.
Posons $B'_\lambda=B_0\otimes_{A_0}A'_\lambda$,
$\mathfrak J'_\lambda=\mathfrak J_0B'_\lambda$,
$M'_\lambda=M_0\otimes_{A_0}A'_\lambda=M_\lambda\otimes_{A_\lambda}A'_\lambda$~;
on a encore $B=\varinjlim B'_\lambda$, $M=\varinjlim M'_\lambda$. Supposons
que l'on ait prouvé qu'il existe un $\lambda$ tel que, pour $\mu\geq\lambda$,
$\operatorname{gr}_{\mathfrak J'_\mu}^*(M'_\mu)$ soit un $A'_\mu$-module
plat. Soit $\mathfrak a_\lambda$ le noyau de l'homomorphisme $\varphi_\lambda$,
qui est donc un idéal de type fini de l'anneau noethérien $A_\lambda$. Il
résulte de la définition de la limite inductive qu'il existe un indice
$\mu\geq\lambda$ tel que $\varphi_{\mu\lambda}(\mathfrak a_\lambda)=0$,
autrement dit l'homomorphisme $\varphi_{\mu\lambda}$ se factorise en
$\varphi_{\mu\lambda}:A_\lambda\to A'_\lambda\to A_\mu$, et l'on peut écrire
$B_\mu=B'_\lambda\otimes_{A'_\lambda}A_\mu$,
$\mathfrak J_\mu=\mathfrak J'_\lambda B_\mu$, et
$M_\mu=M'_\lambda\otimes_{A'_\lambda}A_\mu$~; on déduit donc de (11.2.9.2)
que $\operatorname{gr}_{\mathfrak J_\mu}^*(M_\mu)$ est un $A_\mu$-module
plat.

\medskip
\noindent\textbf{IV)} \emph{Réduction au cas où
$B_0=A_0[T_1,\ldots,T_n]$ et
$\operatorname{gr}_{\mathfrak J_0}^*(B_0)=B_0$.} --- En vertu de l'hypothèse,
il y a un système $(t_i)_{1\leq i\leq n}$ de générateurs de la $A_0$-algèbre
de type fini $B_0$, tel que les $t_i$ pour $i\leq m$ engendrent l'idéal
$\mathfrak J_0$ de $B_0$. Posons $B'_0=A_0[T_1,\ldots,T_n]$,

